% Chapter Template

\chapter{Resultados y Conclusiones} % Main chapter title

%----------------------------------------------------------------------------------------

%----------------------------------------------------------------------------------------
%	SECTION 1
%----------------------------------------------------------------------------------------

\section{Resultados}

\begin{table}[H]
    \centering
    \caption{Métricas de desempeño del modelo en el conjunto de \textit{test}.}
    \label{tab:test_metrics}
    \begin{tabular}{l c}
        \toprule
        \textbf{Métrica} & \textbf{Valor} \\
        \midrule
        Accuracy                & 0.8370103556956326 \\
        F1 Macro                & 0.5564744832188099 \\
        F1 Weighted             & 0.8268692682557662 \\
        Precision Macro         & 0.6166535066964709 \\
        Precision Weighted      & 0.8317142793735933 \\
        Recall Macro            & 0.5467431296475428 \\
        Recall Weighted         & 0.8370103556956326 \\
        \bottomrule
    \end{tabular}
\end{table}

El modelo alcanzó un accuracy de 0.837, demostrando un buen desempeño general. Sin embargo, el F1 Macro = 0.556 indica que el rendimiento no es homogéneo entre clases: algunas categorías, especialmente las menos representadas, presentan resultados más bajos. Por otro lado, el F1 Weighted = 0.827 y las métricas ponderadas de precisión y recall muestran que el modelo funciona muy bien en las clases mayoritarias.

El modelo es sólido para clases frecuentes, pero deficiente en el resto de las clases, esto podria corregirse por lo menos parcialmente con mayor data augmentation, oversampling y un mayor tiempo de entrenamiento.

\section{Conclusiones}

\subsection{Entrenamiento y validación}

\begin{figure}[H]
    \centering
    \includegraphics[width=\textwidth]{figures/loss_curves.png}
    \caption{
        Curvas de \textit{train loss} y \textit{validation loss} para los tres experimentos evaluados:
        \texttt{vit\_articleType\_no\_aug}, \texttt{vit\_aug\_loss} y \texttt{vit\_aug\_v1}.
        Se observa que el modelo sin aumentación presenta sobreajuste temprano, mientras que las variantes con 
        \textit{data augmentation} muestran mejor estabilidad y menor pérdida en validación. 
        Asimismo, \texttt{vit\_aug\_v1} utiliza \textit{accuracy} como métrica de selección del mejor modelo, lo cual no garantiza que el punto guardado corresponda al mínimo de pérdida; 
        por ello no se considera óptimo continuar su entrenamiento bajo ese criterio.
    }
    \label{fig:loss_curves}
\end{figure}

La Figura~\ref{fig:loss_curves} presenta las curvas de \textit{train loss} y \textit{validation loss} para los tres experimentos evaluados: 
\texttt{vit\_articleType\_no\_aug}, \texttt{vit\_aug\_loss} y \texttt{vit\_aug\_v1}. 
La Figura~\ref{fig:val_loss_comparison} complementa este análisis mostrando, en un único gráfico, la evolución conjunta del \textit{val loss} para todas las variantes.

En el caso sin aumentación de datos (\texttt{vit\_articleType\_no\_aug}), el descenso pronunciado del \textit{train loss} acompañado de un \textit{val loss} inestable revela la aparición temprana de sobreajuste y una limitada capacidad de generalización. 
Por el contrario, las variantes con \textit{data augmentation} muestran curvas más estables y valores inferiores de \textit{val loss}, indicando un entrenamiento más robusto ante la variabilidad de los datos.

Entre las configuraciones con aumentación, \texttt{vit\_aug\_loss} selecciona el mejor modelo utilizando \textit{val\_loss}, lo cual constituye el criterio adecuado para asegurar la elección del punto óptimo de entrenamiento. 
En cambio, \texttt{vit\_aug\_v1} utiliza \textit{accuracy} como métrica de selección, que no necesariamente coincide con el mínimo de pérdida y, por tanto, puede llevar a conservar un modelo subóptimo. 
La Figura~\ref{fig:val_loss_comparison} confirma que ambas variantes con augmentations alcanzan pérdidas más bajas que la versión sin aumentación, aunque sólo \texttt{vit\_aug\_loss} garantiza consistencia en la selección del mejor modelo.

\begin{figure}[H]
    \centering
    \includegraphics[width=\textwidth]{figures/loss_curves.png}
    \caption{
        Curvas de entrenamiento y validación para los tres experimentos evaluados. 
        Se observa sobreajuste en el modelo sin augmentations y un comportamiento más estable en las variantes con \textit{data augmentation}. 
        Asimismo, el experimento \texttt{vit\_aug\_v1} utiliza \textit{accuracy} como métrica de selección, lo cual no garantiza la elección del mínimo \textit{val loss}.
    }
    \label{fig:loss_curves}
\end{figure}

\begin{figure}[H]
    \centering
    \includegraphics[width=\textwidth]{figures/val_loss_comparison.png}
    \caption{
        Comparación directa del \textit{validation loss} entre los tres experimentos. 
        Las variantes con augmentations muestran pérdidas sistemáticamente menores, destacándose \texttt{vit\_aug\_loss} por utilizar un criterio de selección coherente con el mínimo de pérdida.
    }
    \label{fig:val_loss_comparison}
\end{figure}

\subsection{Atención}

La Figura~\ref{fig:attention_map} presenta un mapa de atención generado para una de las imágenes del conjunto de prueba. 
En este ejemplo se observa claro que el modelo concentra la mayor parte de su activación sobre los bordes los cuales contienen mucha mas información. 
Asimismo, el fondo es correctamente ignorado, indicando que el modelo aprendió a descartar información irrelevante para la clasificación.

No obstante, también se muestra atención en zonas que no aportan información tan relevante, como partes del cuerpo humano. 
Este comportamiento sugiere que, si bien el mecanismo de atención logra enfocarse en áreas clave, aún distribuye parte de sus recursos en regiones secundarias, lo que podría introducir ruido en la representación interna del modelo. 

\begin{figure}[H]
    \centering
    \includegraphics[width=0.65\textwidth]{figures/attention_map.png}
    \caption{
        Mapa de atención correspondiente a una imagen de ejemplo. 
    }
    \label{fig:attention_map}
\end{figure}

\subsection{Limitaciones y recomendaciones finales}

El modelo presentado podría beneficiarse de un entrenamiento más extenso, ya que en este trabajo nos limitamos a 10 épocas. Un mayor número de iteraciones nos ayudaría a comprobar si la arquitectura continúa mejorando o comienza a sobreajustar. Asimismo, se recomienda comparar el desempeño obtenido con otras arquitecturas de referencia (otros vision transformers, redes convolucionales clásicas u otros modelos híbridos).

Desde una perspectiva de producto, una extensión natural implicaría integrar el modelo en un sistema de entrenamiento continuo, orquestado mediante herramientas de automatización y \textit{job schedulers}. Este \textit{pipeline} permitiría incorporar de forma periódica nuevos datos etiquetados, reentrenar o ajustar el modelo de manera incremental y desplegar versiones actualizadas sin intervención manual, mejorando la capacidad de adaptación del sistema a cambios en la distribución de los datos.

Finalmente, deben señalarse las limitaciones impuestas por el entorno de procesamiento local utilizado. Las restricciones de memoria y capacidad de cómputo condicionaron tanto el tamaño de los lotes como la complejidad de las arquitecturas y la duración de los entrenamientos. El acceso a infraestructuras de cómputo más potentes permitiría experimentos más extensivos y considerables mejoras en el modelo.